% PATCH: C-88
% Reviewer: — (final round-2 corrections; brief intake/pending/final_corrections_round2_2026-11-04.md)
% Target section: sections/results.tex lines 868-1055 (Quantitative Performance Analysis in Physical
%   Multi-Stimulus and Variable Terrain Environments, energy, terrain latency, Limitations)
% Status: applied
%
% Applied, by audit ID (details under "### <ID>" in intake/pending/final_audit_r2_*):
% - B.1 + S-09: Appetitive physical = four replicates (intro sentence, figure caption, table caption);
%   tab:consolidated_physical_metrics Appetitive T_switch cell -> ~0.9-13 s, read from the paper's own
%   ECDF. Aversive, Complex and Obstacle stay at three.
% - S-01 = B-3 = D-05: the 5.06 J figure is presented as a mean, not a worst case or upper bound (no
%   maximum quoted, per the brief); "a typical 38 s segment ... 216 J" -> "the single recorded 38.6 s
%   segment ... ~220 J" (5.71 W x 38.6 s, both from the paper).
% - S-10 = B-1 = D-23: "27x Obstacle baseline" -> "peak" (only that word).
% - B.17c / L-01: "onboard INA sensor" -> "onboard INA219 current sensor" in the two captions the
%   decision names (the table row is added in C-83).
% - S-22: "instantaneous step response" -> "a tightly clustered response (~184 ms)".
% - L-43: the longer DR->Q latency is attributed, not asserted ("which we attribute to").
% - R-24 = S-32: removed the four-panel layout justification from the prose.
% - B.4: lambda in Limitations renamed lambda_conflict (lambda_trial in tab:fsm_comparison, C-86).
% - D-24 / S-37 / L-20: the ablation is referenced through fig:ablation_combined instead of the whole
%   Results section. L-17: "threshold only" without hyphen, as in the rest of Results.
% - R-16: the Limitations closing sentence shortened to a pointer.
% - L-16, L-18, L-31, S-39: Multi-Stimulus, basal-ganglia-inspired, "lateral (roll) axis",
%   "uncontrolled", "rapid backward retreat" (no "computational ticks").
%
% Unified diff (sections/ only):
% diff --git a/sections/results.tex b/sections/results.tex
% index 90cbd51..5919f33 100644
% --- a/sections/results.tex
% +++ b/sections/results.tex
% @@ -860,10 +860,10 @@ To further consolidate the physical validation under complex, multi-section topo
%    \label{fig:real_variable_topography_sequence}
%  \end{figure*}
%  
% -\subsection{Quantitative Performance Analysis in Physical Multi-Stimuli and Variable Terrain Environments}
% +\subsection{Quantitative Performance Analysis in Physical Multi-Stimulus and Variable Terrain Environments}
%  \label{ssec:QuantReal}
%  
% -\revblue{This subsection reports four stimulus scenarios on the physical platform (three hardware replicates each, consolidated in Table~\ref{tab:consolidated_physical_metrics}), followed by energy and latency profiling over variable terrain.}
% +\revblue{This subsection reports four stimulus scenarios on the physical platform (four hardware replicates for Appetitive and three for each of the other scenarios, consolidated in Table~\ref{tab:consolidated_physical_metrics}), followed by energy and latency profiling over variable terrain.}
%  
%  %% ─── APPETITIVE SCENARIO ───────────────────────────────────────────────────
%  
% @@ -882,7 +882,7 @@ To further consolidate the physical validation under complex, multi-section topo
%      \caption{Pitch/Roll RMS attitude dynamics.}
%      \label{fig:appetitive_rms_real}
%    \end{subfigure}
% -  \caption{Real-world Appetitive Stimulus Scenario (three hardware replicates). (a) ECDF of locomotion mode-switching latency $T_{switch}$. (b) Coupled neural and kinematic response: (top) temporal evolution of the basal ganglia firing variance $Var(z)$ with red/blue stimulus-presence shading; (bottom) Roll and Pitch RMS attitude profile throughout the trial.}
% +  \caption{Real-world Appetitive Stimulus Scenario (\revblue{four} hardware replicates). (a) ECDF of locomotion mode-switching latency $T_{switch}$. (b) Coupled neural and kinematic response: (top) temporal evolution of the basal ganglia firing variance $Var(z)$ with red/blue stimulus-presence shading; (bottom) Roll and Pitch RMS attitude profile throughout the trial.}
%    \label{fig:appetitive_real_compound}
%  \end{figure*}
%  
% @@ -909,7 +909,7 @@ To further consolidate the physical validation under complex, multi-section topo
%    \label{fig:aversive_real_compound}
%  \end{figure*}
%  
% -\textbf{Aversive (see Figure~\ref{fig:aversive_real_compound}):} Unlike the Appetitive scenario's extended settling window, the Aversive pipeline behaves as an immediate emergency reflex: upon threat detection, the ROS~2 Control architecture bypasses heavy trajectory filtering to prioritize high-speed backward flight, minimizing routing overhead to near-zero computational ticks under critical hazards. The dual-peak firing-variance pattern reflects, first, detection of the aversive stimulus and dispatch of the mode-switching command, then a second, higher-energy peak as hardware inertia keeps the threat briefly within the camera's field of view during the retreat, decaying back to its resting state once the chassis clears the threat from view. The concurrent Pitch step-increase and cyclic Roll oscillations reflect transient longitudinal tilt during the wheel-driven evasion, returning to a flat baseline once the threat leaves the visual field.
% +\textbf{Aversive (see Figure~\ref{fig:aversive_real_compound}):} Unlike the Appetitive scenario's extended settling window, the Aversive pipeline behaves as an immediate emergency reflex: upon threat detection, the ROS~2 Control architecture bypasses heavy trajectory filtering to prioritize a rapid backward retreat, minimizing routing overhead under critical hazards. The dual-peak firing-variance pattern reflects, first, detection of the aversive stimulus and dispatch of the mode-switching command, then a second, higher-energy peak as hardware inertia keeps the threat briefly within the camera's field of view during the retreat, decaying back to its resting state once the chassis clears the threat from view. The concurrent Pitch step-increase and cyclic Roll oscillations reflect transient longitudinal tilt during the wheel-driven evasion, returning to a flat baseline once the threat leaves the visual field.
%  
%  %% ─── COMPLEX MULTI-STIMULUS SCENARIO ─────────────────────────────────────
%  
% @@ -932,7 +932,7 @@ To further consolidate the physical validation under complex, multi-section topo
%    \label{fig:complex_real_compound}
%  \end{figure*}
%  
% -\textbf{Complex (see Figure~\ref{fig:complex_real_compound}):} \revblue{The step response stays deterministic across back-to-back transformations: the ROS~2 Control layer absorbs the electromechanical inertia and blends the trajectories, so torque stress does not accumulate in the joints.} The three-phase firing-variance profile maps the multi-stage decision process: a transient hyper-excitation peak during static obstacle avoidance, \revblue{driven by the dense LiDAR returns from the physical barrier}; a secondary activation wave capturing the emergency flight reflex to the aversive cue; and a low-energy steady-state as the robot tracks the appetitive target to mission completion. The chassis attitude telemetry mirrors this sequence, with asymmetric flanking maneuvers exciting the transversal axis and rapid backward accelerations peaking the longitudinal tilt, before both metrics drop and stabilize during the final convergence phase, indicating the wheel-driven configuration smooths parasitic oscillations during target approach.
% +\textbf{Complex (see Figure~\ref{fig:complex_real_compound}):} \revblue{The step response stays deterministic across back-to-back transformations: the ROS~2 Control layer absorbs the electromechanical inertia and blends the trajectories, so torque stress does not accumulate in the joints.} The three-phase firing-variance profile maps the multi-stage decision process: a transient hyper-excitation peak during static obstacle avoidance, \revblue{driven by the dense LiDAR returns from the physical barrier}; a secondary activation wave capturing the emergency flight reflex to the aversive cue; and a low-energy steady-state as the robot tracks the appetitive target to mission completion. The chassis attitude telemetry mirrors this sequence, with asymmetric flanking maneuvers exciting the lateral (roll) axis and rapid backward accelerations peaking the longitudinal tilt, before both metrics drop and stabilize during the final convergence phase, indicating the wheel-driven configuration smooths parasitic oscillations during target approach.
%  
%  %% ─── STATIC OBSTACLE AVOIDANCE SCENARIO ──────────────────────────────────
%  
% @@ -959,15 +959,15 @@ To further consolidate the physical validation under complex, multi-section topo
%  
%  \begin{table}[htbp]
%  \centering
% -\caption{Consolidated Quantitative Metrics for the Physical Multi-Stimuli Scenarios (three hardware replicates each). $T_{switch}$: mode-switching-delay ECDF window. $Var(z)$: basal-ganglia firing-variance peak/pattern. Roll/Pitch RMS: chassis attitude range during the active response window.}
% +\caption{Consolidated Quantitative Metrics for the Physical Multi-Stimulus Scenarios (\revblue{four hardware replicates for Appetitive, three for each other scenario}). $T_{switch}$: mode-switching-delay ECDF window. $Var(z)$: basal-ganglia firing-variance peak/pattern. Roll/Pitch RMS: chassis attitude range during the active response window.}
%  \label{tab:consolidated_physical_metrics}
%  \begin{tabularx}{\textwidth}{l >{\centering\arraybackslash}X >{\centering\arraybackslash}X >{\centering\arraybackslash}X >{\centering\arraybackslash}X}
%  \toprule
%  \textbf{Scenario} & \textbf{$T_{switch}$ ECDF} & \textbf{$Var(z)$ peak} & \textbf{Roll / Pitch RMS} & \textbf{Settles by} \\
%  \midrule
% -Appetitive & $3.0$--$9.0\,\text{s}$ (mechanically limited) & Sustained, non-zero post-completion & $3.0^{\circ}$--$3.5^{\circ}$ (cyclic) & $t > 25\,\text{s}$ \\
% +Appetitive & \revblue{${\approx}\,0.9$--$13\,\text{s}$} (mechanically limited) & Sustained, non-zero post-completion & $3.0^{\circ}$--$3.5^{\circ}$ (cyclic) & $t > 25\,\text{s}$ \\
%  Aversive   & ${\approx}\,0.0\,\text{s}$ (instantaneous) & Dual-peak (detection + retreat-FOV) & Pitch step $\uparrow$, Roll cyclic ($t{=}4$--$9\,\text{s}$) & $t > 10\,\text{s}$ \\
% -Complex    & ${\approx}\,1.0$--$1.2\,\text{s}$ & Three-phase, peak ${\approx}\,27\times$ Obstacle baseline & Roll ${\approx}\,7^{\circ}$, Pitch ${\approx}\,5^{\circ} \to 4^{\circ}$ & $t > 25\,\text{s}$ \\
% +Complex    & ${\approx}\,1.0$--$1.2\,\text{s}$ & Three-phase, peak ${\approx}\,27\times$ Obstacle \revblue{peak} & Roll ${\approx}\,7^{\circ}$, Pitch ${\approx}\,5^{\circ} \to 4^{\circ}$ & $t > 25\,\text{s}$ \\
%  Obstacle   & ${\approx}\,0.0\,\text{s}$ (instantaneous) & ${\approx}\,1100$, transient ($t < 7\,\text{s}$) & Roll $3.5^{\circ}$ (rhythmic) & $t = 8\,\text{s}$ \\
%  \bottomrule
%  \end{tabularx}
% @@ -978,18 +978,18 @@ Obstacle   & ${\approx}\,0.0\,\text{s}$ (instantaneous) & ${\approx}\,1100$, tra
%  \begin{figure}[t]
%    \centering
%    \includegraphics[width=0.5\textwidth]{FigReal_Morphological_Transition_Energy.png}
% -  \caption{Morphological transition energy per direction, derived from $|I_{\text{current}}| \times 7.2\,\text{V}$ \revblue{(current sampled at $\approx 5$\,Hz by the onboard INA sensor; fixed nominal voltage of the NiMH battery, Table~\ref{table:HardwareParams})} integrated over the pre-calibrated electromechanical coupling window (2.92\,s / 1.56\,s, matching Figure~\ref{fig:FourPanelLatencies}), across $N=16$ (Quadrupedal $\rightarrow$ Differential Rolling) and $N=19$ (Differential Rolling $\rightarrow$ Quadrupedal) individually detected structural-command transitions. Error bars: $\pm 1$ SD; the between-direction difference is not statistically significant (Welch's $t \approx 1.18$).}
% +  \caption{Morphological transition energy per direction, derived from $|I_{\text{current}}| \times 7.2\,\text{V}$ \revblue{(current sampled at $\approx 5$\,Hz by the onboard INA219 current sensor; fixed nominal voltage of the NiMH battery, Table~\ref{table:HardwareParams})} integrated over the pre-calibrated electromechanical coupling window (2.92\,s / 1.56\,s, matching Figure~\ref{fig:FourPanelLatencies}), across $N=16$ (Quadrupedal $\rightarrow$ Differential Rolling) and $N=19$ (Differential Rolling $\rightarrow$ Quadrupedal) individually detected structural-command transitions. Error bars: $\pm 1$ SD; the between-direction difference is not statistically significant (Welch's $t \approx 1.18$).}
%    \label{fig:MorphologicalEnergy}
%  \end{figure}
%  
% -The energetic cost of physical adaptation was quantified from individually detected structural-command transitions \revblue{as described in Figure~\ref{fig:MorphologicalEnergy}'s caption}. The Quadrupedal $\rightarrow$ Differential Rolling transition required $5.06 \pm 2.93$~J ($N=16$), and the reverse transition (Differential Rolling $\rightarrow$ Quadrupedal) required $4.13 \pm 1.19$~J ($N=19$); the between-direction difference is not statistically significant \revblue{(see caption)}. Greater trial-to-trial variability was observed in the Quadrupedal $\rightarrow$ Differential Rolling direction, though the present dataset does not isolate a specific mechanical cause for this dispersion. As the Quadrupedal$\leftrightarrow$Differential Rolling pair involves the largest morphological reconfiguration among the three modes (full limb deployment/retraction together with the leg-to-wheel locomotion change), its transition energy serves as an upper bound on the cost of any mode transition: in the worst case a transition consumes $\leq 5.06$~J, a small cost relative to continuous operation (e.g.\ a typical 38\,s Omnidirectional segment integrates to $\approx 216$~J). \revblue{Table~\ref{tab:mode_steady_state_energy} reports the long-term average steady-state power draw of each mode.}
% +The energetic cost of physical adaptation was quantified from individually detected structural-command transitions \revblue{as described in Figure~\ref{fig:MorphologicalEnergy}'s caption}. The Quadrupedal $\rightarrow$ Differential Rolling transition required $5.06 \pm 2.93$~J ($N=16$), and the reverse transition (Differential Rolling $\rightarrow$ Quadrupedal) required $4.13 \pm 1.19$~J ($N=19$); the between-direction difference is not statistically significant \revblue{(see caption)}. Greater trial-to-trial variability was observed in the Quadrupedal $\rightarrow$ Differential Rolling direction, though the present dataset does not isolate a specific mechanical cause for this dispersion. \revblue{The Quadrupedal$\leftrightarrow$Differential Rolling pair involves the largest morphological reconfiguration among the three modes (full limb deployment/retraction together with the leg-to-wheel locomotion change); its mean transition energy, 5.06~J in the more costly direction, is small relative to continuous operation (e.g.\ the single recorded 38.6\,s Omnidirectional segment integrates to $\approx 220$~J).} \revblue{Table~\ref{tab:mode_steady_state_energy} reports the long-term average steady-state power draw of each mode.}
%  
%  \begin{table}[htbp]
%  \centering
%  \caption{Long-term average steady-state power consumption per locomotion mode, computed from
%  pooled intra-segment power samples (individual samples within each structural-command-detected mode
%  segment, electromechanical coupling window excluded; see Figure~\ref{fig:MorphologicalEnergy}). Power
% -derived as $|I|\times7.2\,\text{V}$\revblue{, with current sampled at $\approx 5$\,Hz by the onboard INA sensor and 7.2\,V the nominal voltage of the 3300\,mAh NiMH battery (Table~\ref{table:HardwareParams}). Mean and SD are taken over all individual samples pooled across the segments of each mode.}}
% +derived as $|I|\times7.2\,\text{V}$\revblue{, with current sampled at $\approx 5$\,Hz by the onboard INA219 current sensor and 7.2\,V the nominal voltage of the 3300\,mAh NiMH battery (Table~\ref{table:HardwareParams}). Mean and SD are taken over all individual samples pooled across the segments of each mode.}}
%  \label{tab:mode_steady_state_energy}
%  \begin{tabularx}{\textwidth}{l >{\centering\arraybackslash}X >{\centering\arraybackslash}X}
%  \toprule
% @@ -1017,12 +1017,12 @@ This roughly five-fold gap between Omnidirectional and Differential Rolling has
%    \label{fig:FourPanelLatencies}
%  \end{figure}
%  
% -To evaluate repeatability and determinism under uncalibrated physical environments, a continuous multi-terrain navigation mission was executed across 15 hardware replicates. The track forced the robot from Quadrupedal to Differential Rolling Mode to climb a steep incline, then back to Quadrupedal Mode on unstructured terrain inducing wheel slippage; the resulting temporal dynamics are partitioned into a four-panel layout to preserve resolution against low statistical variance (Figure~\ref{fig:FourPanelLatencies}). The basal ganglia loop's decision latency shows a clear task-dependent split: the Quadrupedal $\rightarrow$ Differential Rolling transition (Panel~A) achieves an instantaneous step response, reflecting immediate pitch-gradient adaptation, while the reverse transition (Panel~B), triggered by wheel slippage on rugged terrain, settles at a markedly longer latency, a deliberate stability feature that filters transient high-frequency impact noise to suppress control chattering and false-positive switches on noisy surfaces. The electromechanical coupling timelines reflect these commands mechanically: achieving full wheel-ground alignment (Panel~C) requires an extended settling window, capturing the kinematic cost of leg extension under slope conditions, whereas the emergency retraction to the high-traction articulated stance (Panel~D) is a rapid recovery reflex, minimizing the wheel-spin window and securing foothold anchors quickly on unstable terrain.
% +To evaluate repeatability and determinism under uncalibrated physical environments, a continuous multi-terrain navigation mission was executed across 15 hardware replicates. The track forced the robot from Quadrupedal to Differential Rolling Mode to climb a steep incline, then back to Quadrupedal Mode on unstructured terrain inducing wheel slippage; the resulting latencies are shown in Figure~\ref{fig:FourPanelLatencies}. The basal ganglia loop's decision latency shows a clear task-dependent split: the Quadrupedal $\rightarrow$ Differential Rolling transition (Panel~A) achieves \revblue{a tightly clustered response ($\approx 184$~ms)}, reflecting immediate pitch-gradient adaptation, while the reverse transition (Panel~B), triggered by wheel slippage on rugged terrain, settles at a markedly longer latency, \revblue{which we attribute to} a deliberate stability feature that filters transient high-frequency impact noise to suppress control chattering and false-positive switches on noisy surfaces. The electromechanical coupling timelines reflect these commands mechanically: achieving full wheel-ground alignment (Panel~C) requires an extended settling window, capturing the kinematic cost of leg extension under slope conditions, whereas the emergency retraction to the high-traction articulated stance (Panel~D) is a rapid recovery reflex, minimizing the wheel-spin window and securing foothold anchors quickly on unstable terrain.
%  
%  \subsection{Limitations of the Present Work}
%  \label{ssec:Limitations}
%  
% -As an open, low-cost exploratory research baseline, the physical robot is constructed using low-torque servomotors and a lightweight PLA structure, which limits its operational envelope compared to industrial-grade platforms. Locomotion gaits are currently executed via pre-configured open-loop control patterns; while this reduces control complexity, it weakens active stabilization capabilities on high-slip surfaces or excessively steep inclinations. \revblue{This limit lies in the hardware and is independent of the neural arbitration architecture.} As a representative case, the gait selected during steep-slope negotiation (Figure~\ref{fig:RES_IMU2_real}) reflects the boundaries of the physical actuator envelope: the neural decision module consistently selects the locomotion mode associated with improved stability under the detected terrain conditions, yet the open-loop execution of that gait cannot compensate for the torque deficits inherent to the hardware. Experimental validation was constrained to structured laboratory environments with static stimuli and did not include unconstrained outdoor topographies, dynamic obstacles, or severe external disturbances. Therefore, the empirical results presented herein do not constitute validation of a production-ready system for rugged industrial deployment. Instead, they substantiate the mathematical, biological, and computational feasibility of a highly efficient basal-ganglia inspired arbitration network capable of managing multimodal mobility \revblue{on a low-cost robotic platform}.
% +As an open, low-cost exploratory research baseline, the physical robot is constructed using low-torque servomotors and a lightweight PLA structure, which limits its operational envelope compared to industrial-grade platforms. Locomotion gaits are currently executed via pre-configured open-loop control patterns; while this reduces control complexity, it weakens active stabilization capabilities on high-slip surfaces or excessively steep inclinations. \revblue{This limit lies in the hardware and is independent of the neural arbitration architecture.} As a representative case, the gait selected during steep-slope negotiation (Figure~\ref{fig:RES_IMU2_real}) reflects the boundaries of the physical actuator envelope: the neural decision module consistently selects the locomotion mode associated with improved stability under the detected terrain conditions, yet the open-loop execution of that gait cannot compensate for the torque deficits inherent to the hardware. Experimental validation was constrained to structured laboratory environments with static stimuli and did not include unconstrained outdoor topographies, dynamic obstacles, or severe external disturbances. Therefore, the empirical results presented herein do not constitute validation of a production-ready system for rugged industrial deployment. Instead, they substantiate the mathematical, biological, and computational feasibility of a highly efficient basal-ganglia-inspired arbitration network capable of managing multimodal mobility \revblue{on a low-cost robotic platform}.
%  
%  
%  \revblue{Regarding the computing platform: the neural circuits, the MLP terrain classifier, and the image processing ran on the external laptop of Section~\ref{ssec:HardwareSpec}, not on the robot's ESP32-S3. Onboard feasibility was assessed only for the MLP classifier (Section~\ref{sssec:mlp_eval}); the basal-ganglia, locomotion, and gait-decision circuits and the vision pipeline were not benchmarked on the microcontroller, so a fully autonomous onboard implementation remains future work.}
% @@ -1042,9 +1042,9 @@ constrained environment, rather than as performance guarantees
%  for deployment in specific industrial buildings.%
%  
%  
% -Beyond the hardware and simulation-environment constraints above, the neural and locomotion-control parameters were manually tuned and kept constant throughout all experiments, without applying adaptive optimization processes that could enhance performance in more diverse scenarios; system resilience was also not assessed under partial hardware failures or temporary sensory dropout conditions likely to occur in real industrial environments \revblue{(beyond the simulated camera-detection dropout of Section~\ref{sssec:noise_robustness_extra})}. The multimodal sensory integration, moreover, does not follow a formal fusion framework in the classical signal-processing sense, such as Kalman filtering, Bayesian estimation, or deep sensor fusion. Instead, cross-modal integration emerges implicitly within the bio-inspired neural network: each sensory channel feeds differentiated neural pathways that converge through winner-take-all competition dynamics in the basal ganglia module, effectively distributing the fusion function across the network topology. While this approach is biologically grounded, it has not been formally compared against standard fusion methods. Similarly, visual stimulus detection relies on color-based segmentation calibrated under controlled illumination, without evaluation under varying lighting conditions or occlusion. Notwithstanding these design simplifications, the perception-decision noise-robustness evaluation (Section~\ref{ssec:NoiseRobustnessEval}) demonstrated empirical robustness under combined IMU-drift, LiDAR-noise, and camera-illumination-distortion perturbation, \revblue{applied jointly: no configuration isolates a single sensor's degradation.} Within the same evaluation, additive IMU white noise (distinct from the random-walk drift that was exercised) was held at zero throughout and so was not tested here. Robustness was also not uniform across scenarios: the Rugged Terrain suite showed a markedly elevated failure rate under perturbation relative to the other five suites, with 8 tipovers and 4 crashes recorded out of 27 attempts (Section~\ref{sssec:noise_robustness_extra}), indicating that combined sensory perturbation degrades structural stability specifically on this terrain even though the trials that do complete show no corresponding trend toward instability. These constraints do not invalidate the results but indicate that the next development stage should focus on validating the proposed architecture in a prototype that includes robustness and adaptability tests under real, non-controlled conditions.
% +Beyond the hardware and simulation-environment constraints above, the neural and locomotion-control parameters were manually tuned and kept constant throughout all experiments, without applying adaptive optimization processes that could enhance performance in more diverse scenarios; system resilience was also not assessed under partial hardware failures or temporary sensory dropout conditions likely to occur in real industrial environments \revblue{(beyond the simulated camera-detection dropout of Section~\ref{sssec:noise_robustness_extra})}. The multimodal sensory integration, moreover, does not follow a formal fusion framework in the classical signal-processing sense, such as Kalman filtering, Bayesian estimation, or deep sensor fusion. Instead, cross-modal integration emerges implicitly within the bio-inspired neural network: each sensory channel feeds differentiated neural pathways that converge through winner-take-all competition dynamics in the basal ganglia module, effectively distributing the fusion function across the network topology. While this approach is biologically grounded, it has not been formally compared against standard fusion methods. Similarly, visual stimulus detection relies on color-based segmentation calibrated under controlled illumination, without evaluation under varying lighting conditions or occlusion. Notwithstanding these design simplifications, the perception-decision noise-robustness evaluation (Section~\ref{ssec:NoiseRobustnessEval}) demonstrated empirical robustness under combined IMU-drift, LiDAR-noise, and camera-illumination-distortion perturbation, \revblue{applied jointly: no configuration isolates a single sensor's degradation.} Within the same evaluation, additive IMU white noise (distinct from the random-walk drift that was exercised) was held at zero throughout and so was not tested here. Robustness was also not uniform across scenarios: the Rugged Terrain suite showed a markedly elevated failure rate under perturbation relative to the other five suites, with 8 tipovers and 4 crashes recorded out of 27 attempts (Section~\ref{sssec:noise_robustness_extra}), indicating that combined sensory perturbation degrades structural stability specifically on this terrain even though the trials that do complete show no corresponding trend toward instability. These constraints do not invalidate the results but indicate that the next development stage should focus on validating the proposed architecture in a prototype that includes robustness and adaptability tests under real, uncontrolled conditions.
%  
% -Regarding the basal-ganglia ablation study (Section~\ref{sec:results}): the conflict-efficiency metric $\lambda$ was excluded from the reported results because its current implementation only registers active conflict when two specific channels are simultaneously present within the sampling window, a condition rarely met in practice across these trials; this is a limitation of the metric's present instrumentation, not evidence that conflict efficiency is unaffected by the ablations. Additionally, per-cycle firing-variance and temporal-consistency diagnostics, which could further characterize the \emph{threshold-only} variant's failure mode (e.g.\ oscillation near the decision boundary, as opposed to outright deadlock), were not captured for this dataset; the threshold-versus-WTA comparison therefore relies on decision latency and attitude stability rather than on these more targeted diagnostics.
% +Regarding the basal-ganglia ablation study (Figure~\ref{fig:ablation_combined}): the conflict-efficiency metric \revblue{$\lambda_{conflict}$} was excluded from the reported results because its current implementation only registers active conflict when two specific channels are simultaneously present within the sampling window, a condition rarely met in practice across these trials; this is a limitation of the metric's present instrumentation, not evidence that conflict efficiency is unaffected by the ablations. Additionally, per-cycle firing-variance and temporal-consistency diagnostics, which could further characterize the \emph{threshold only} variant's failure mode (e.g.\ oscillation near the decision boundary, as opposed to outright deadlock), were not captured for this dataset; the threshold-versus-WTA comparison therefore relies on decision latency and attitude stability rather than on these more targeted diagnostics.
%  
% -The Reactive Inspection Task's scope constraints (absence of deliberate path planning or mapping, the resulting inapplicability of a formal coverage-rate metric, and the first- rather than re-acquisition semantics of $T_{acquisition}$) are established alongside that experiment's own results in Section~\ref{sssec:reactive_inspection} above and are not restated here.
% +\revblue{The Reactive Inspection Task's scope constraints are stated with its results (Section~\ref{sssec:reactive_inspection}).}
%  
